\documentclass[10pt,a4paper]{amsart}
\usepackage[margin=19mm]{geometry}
\usepackage{amsmath,amssymb,mathtools}
\usepackage[scr=rsfs,cal=euler]{mathalfa}
\usepackage{xfrac}
\usepackage[unicode,hidelinks,bookmarks=false]{hyperref}
\newcommand{\ie}{i.e.\ }
\newcommand{\cf}{cf.\ }
\newcommand{\resp}{respectively\ }
\newcommand{\Subsection}{\subsection}
\providecommand{\nobreakdash}{\nobreak\hbox{-}}
\setlength{\emergencystretch}{2em}
\multlinegap0pt
\usepackage{bm}
\usepackage{bbm}
\usepackage{dsfont}
\usepackage{xspace}
\usepackage{cancel}
\usepackage{mathtools}
\usepackage{amsthm}
\makeatletter
\@ifundefined{theorem}{\newtheorem{theorem}{Theorem}}{}
\makeatother
\title[Rigidity of solitons of the Gauss curvature flow in Euclidea]{Rigidity of solitons of the Gauss curvature flow in Euclidean space}
\author{Rafael L\'opez}
\date{Source: arXiv:2507.12863v1 (CC BY 4.0)}
\begin{document}
\maketitle

\noindent\emph{Source and licence.} Rigidity of solitons of the Gauss curvature flow in Euclidean space. Version arXiv:2507.12863v1 (CC BY 4.0), distributed under the Creative Commons Attribution 4.0 International licence, as recorded in the arXiv licence metadata for this exact version and shown on the abstract page https://arxiv.org/abs/2507.12863v1. Full rights notice: licenses/arxiv-2507.12863-CC-BY-4.0.txt.\par\smallskip

\noindent\textit{Editorial scope.} Counted unit: the Theorem whose source label is \texttt{t-d} in the article (source0.tex lines 464--465, source label \texttt{t-d}), delivered together with the article's own complete original proof of it (source0.tex lines 467--510, one \texttt{proof} environment of 44 lines). No mathematics is written, repaired or re-derived: statement and proof are the article's own text line for line. Required context retained uncounted: b1 (lines 165--172); bb (lines 174--177); d3 (lines 428--438). Every definition, display and label of those lines is retained unchanged. Omitted from this package and not redistributed: 1--164, 173--173, 178--427, 439--463, 466--466, 511--659. Nothing inside a retained excerpt is omitted or altered: every retained body is delivered exactly as the article's lines are written, so no omission receipt of any kind accompanies this package. Citations: the retained excerpts carry no citation command at all, so no bibliography entry of the article is transcribed and no reference is redistributed. Reference adaptation: none was needed and none was made; every reference inside the delivered unit resolves inside it, so no unresolved reference or citation is produced. Printed slips kept literal: no printed word, symbol, hypothesis or number of the counted statement or of its proof was altered. Layout and class: the article's own class is \texttt{amsart} with packages including amssymb, amsthm, amscd, graphicx, color, url; the wrapper is the lane's pinned \texttt{amsart} editorial wrapper at 10pt on A4, which restates the theorem-like environments the retained text needs and adds the article's own macro definitions that the retained text needs (none), with extra wrapper packages (bm, bbm, dsfont, xspace, cancel, mathtools). The delivered numbering is the unit's own. Assets: no retained span contains a figure environment, a graphics-inclusion command, a PDF-page inclusion, a table or a TikZ picture; no image file is copied into this package, the package declares no asset dependency and no image file is redistributed with it. Licence: version arXiv:2507.12863v1 of this article is distributed under the Creative Commons Attribution 4.0 International licence, as recorded in the arXiv licence metadata for this exact version and shown on the abstract page https://arxiv.org/abs/2507.12863v1. A full rights notice is delivered at licenses/arxiv-2507.12863-CC-BY-4.0.txt. Characters: every retained excerpt is pure ASCII, so the delivered engine, XeLaTeX with its default Latin Modern text font, reports no missing character.
\begin{equation}\label{b1}
    \left\{
    \begin{array}{ll}
        e_{2}(\kappa_1)=(\kappa_1-\kappa_2)\omega(e_1)\\
        e_{1}(\kappa_2)=(\kappa_1-\kappa_2)\omega(e_2).
    \end{array}
    \right.
\end{equation}

\begin{equation}\label{bb}
   K=\kappa_1\kappa_2=-e_1\left(\frac{e_1( \kappa_2 )}{\kappa_1-\kappa_2} \right)+e_2\left(\frac{e_2( \kappa_1 )}{\kappa_1 -\kappa_2}\right)
   -\frac{e_1( \kappa_2)^2+ e_2( \kappa_1 )^2}{(\kappa_1 -\kappa_2)^2}.
\end{equation}

    \begin{equation}\label{d3}
    \left\{
\begin{split}
            &e_{1}(\gamma)-\mu\omega(e_1)-\frac{K-\lambda}{\alpha}\kappa_1=1\\
            &e_{2}(\gamma)-\mu\omega(e_2)=0\\
            &e_{1}(\mu)+\gamma\omega(e_1)=0\\
            &e_{2}(\mu)+\gamma\omega(e_2)-\frac{K-\lambda}{\alpha}\kappa_2=1\\
            &e_{1}(\frac{K}{\alpha})+\gamma\kappa_1=0\\
            &e_{2}(\frac{K}{\alpha})+\mu\kappa_2=0.
        \end{split}\right.
    \end{equation}

\begin{theorem}\label{t-d} Let $\Sigma$ be a $\lambda$-shrinker of the GCF. If a principal curvature is constant, then $\Sigma$ is a plane, a sphere or a circular cylinder. 
\end{theorem}

\begin{proof} Without loss of generality, assume that  the principal curvature $\kappa_2$ is constant. Again, if   $\kappa_1$ is constant, then the surface is isoparametric and because $\Sigma$ is a $\lambda$-shrinker, then  $\Sigma$ is a plane, a sphere or a circular cylinder. This proves the result  in this situation.

Suppose now that $\kappa_1$ is not constant and we will arrive to a contradiction.

\begin{enumerate}
\item Case $\kappa_2=0$. Then $K=0$. From \eqref{b1} and   \eqref{d3}, we obtain $\gamma=0$, $\omega(e_2)=0$ and $e_1(\mu)=0$, $e_2(\mu)=1$.      In particular,  $\mu\not=0$ identically. The first equation of \eqref{d3} is 
$\mu\omega(e_1) +\frac{\lambda}{\alpha}\kappa_1=1$. From \eqref{b1} we know $\omega(e_1)=\frac{e_2(\kappa_1)}{\kappa_1}$, hence
\begin{equation}\label{con}
\frac{e_2(\kappa_1)}{\kappa_1}\mu+\frac{\lambda}{\alpha}\kappa_1=1.
\end{equation}
Differentiating with respect to $e_2$ and taking into account that $e_2(\mu)=1$, we have  
\begin{equation}\label{con2}
\frac{e_{22}(\kappa_1)}{\kappa_1}\mu+\frac{e_2(\kappa_1)}{\kappa_1}-\frac{e_2(\kappa_1)^2}{\kappa_1^2}\mu+\frac{\lambda}{\alpha}e_2(\kappa_1)=0.
\end{equation}
On the other hand, the identity \eqref{bb} is 
$$0=\frac{e_{22}(\kappa_1)}{\kappa_1}-2\frac{e_2(\kappa_1)^2}{\kappa_1^2}=0.$$ 
From this equation, we get the value of $e_{22}(\kappa_1)$ and putting in \eqref{con2}, we obtain 
$$\frac{e_2(\kappa_1)^2}{\kappa_1^2}\mu+\frac{e_2(\kappa_1)}{\kappa_1}+\frac{\lambda}{\alpha}e_2(\kappa_1)=0.$$
Because $e_2(\kappa_1)\not=0$, then
$$\frac{e_2(\kappa_1)}{\kappa_1^2}\mu+\frac{1}{\kappa_1}+\frac{\lambda}{\alpha}=0.$$
This gives a contradiction with \eqref{con}.
\item Case $\kappa_2\not=0$.    Let us work on the  open set $\Omega\subset \Sigma$ formed by non-umbilic points. From \eqref{b1} we have $\omega(e_2)=0$. By using the forth and sixth equations of \eqref{d3}, we deduce   $\mu=- e_2(\kappa_1)/\alpha$ and $e_2(\mu)=1+\frac{K-\lambda}{\alpha}\kappa_2$. From the first identity, we have  
$e_2(\mu)=-  e_{22}(\kappa_1)/\alpha$. Then 
\begin{equation}\label{d22}
e_{22}(\kappa_1)=-(K-\lambda)\kappa_2-\alpha.
\end{equation}
Using this identity, and because $\kappa_2$ is constant, equation \eqref{bb} is 
$$\kappa_1\kappa_2=\frac{e_{22}(\kappa_1)}{\kappa_1-\kappa_2}-2\frac{e_2(\kappa_1)^2}{(\kappa_1-\kappa_2)^2}=
-\frac{(K-\lambda)\kappa_2+\alpha}{\kappa_1-\kappa_2}-2\frac{e_2(\kappa_1)^2}{(\kappa_1-\kappa_2)^2}.$$
 Then 
 $$2e_2(\kappa_1)^2=-(\kappa_1-\kappa_2)(\kappa_1 K+\alpha-\lambda\kappa_2).$$
 Differentiating with respect to $e_2$, we have 
 \begin{equation}\label{fp22}
4e_2(\kappa_1)e_{22}(\kappa_1)=-2e_2(\kappa_1)\left(   \alpha -4 H^2 \kappa_1+9 H \kappa_1^2-3 H \lambda -4 \kappa_1^3+2 \lambda  \kappa_1 \right).
\end{equation}
 The function  $e_2(\kappa_1)$ is not zero because otherwise, we have $\mu=0$ and the forth equation in \eqref{d3} yields $K-\lambda=1$. Then  $K$ is constant and  both principal curvatures are constant because $\kappa_2\not=0$. This is a contradiction because this situation was initiated discarded.
 
   
  
 We insert  \eqref{d22} in \eqref{fp22}. Using that $\kappa_2\not=0$, we obtain  a polynomial on the function $\kappa_1$, namely, 
 $$-6 \kappa_1^3+17H\kappa_1^2-12H^2\kappa_1+\lambda H-\alpha =0.$$
  This gives a contradiction because $\kappa_1$ is a non-constant function. We conclude that the case $\kappa_2\not=0$ cannot occur.
 \end{enumerate}
  \end{proof}

\end{document}
